\textbf{Audiencia:} profesorado de profesional y posgrado de la \textbf{EIC}.\\
\textbf{Modalidad:} virtual (Teams), híbrida o presencial según edición.\\
Leyenda (aula): \textbf{E} = exposición / demo · \textbf{P} = práctica · \textbf{C} = cierre / puesta en común

\section{Parte A --- Sesiones en aula (8 $\times$ 3\,h = 24\,h)}

\subsection{Sesión 1 --- Fundamentos agénticos e integridad académica (EIC)}
\textbf{Duración:} 3:00 · \textbf{Producto (E1):} Mapa agentes vs.\ chat + nota de integridad + tema EIC

\begin{tabularx}{\textwidth}{@{}lllX@{}}
\toprule
Min & Bloque & Tipo & Actividad \\
\midrule
0:00--0:15 & Bienvenida y mapa del CADi (40\,h) & E & Objetivos, 24+16\,h, evidencias, reglas EIC \\
0:15--0:45 & Agentes vs.\ chat & E & Roles, herramientas, memoria \\
0:45--1:10 & Integridad académica y riesgos & E & Alucinación, sesgo, datos; ecuaciones, citas \\
1:10--1:20 & \emph{Descanso} & --- & --- \\
1:20--1:45 & Demo: primer flujo con HITL & E & Idea $\rightarrow$ agente $\rightarrow$ humano \\
1:45--2:35 & Taller: mapa + nota de integridad & P & Tema propio EIC \\
2:35--2:50 & Puesta en común & C & 3--4 ejemplos \\
2:50--3:00 & Cierre y TI hacia sesión 2 & C & Lecturas bloque A + bosquejo de brief \\
\bottomrule
\end{tabularx}

\textbf{TI sugerido post-sesión:} $\sim$1\,h lecturas/demos (bloque A)

\subsection{Sesión 2 --- Briefs, prompts de sistema y orquestación}
\textbf{Duración:} 3:00 · \textbf{Producto (E2):} Brief + prompt v0 + esquema de orquestación

\begin{tabularx}{\textwidth}{@{}lllX@{}}
\toprule
Min & Bloque & Tipo & Actividad \\
\midrule
0:00--0:10 & Recap integridad y temas EIC & E & Conexión con E1 \\
0:10--0:40 & Brief como contrato pedagógico & E & Audiencia, alcance, restricciones \\
0:40--1:05 & Prompt de sistema y multi-rol & E & Anatomía; planificador / redactor / revisor \\
1:05--1:15 & \emph{Descanso} & --- & --- \\
1:15--1:40 & Demo: orquestación con HITL & E & Secuencial + crítico-humano \\
1:40--2:30 & Taller: brief + prompt + diagrama & P & Plantillas; 2--4 roles \\
2:30--2:50 & Puesta en común & C & Restricciones y no-citas inventadas \\
2:50--3:00 & Cierre & C & Practicar plantillas (bloque B) \\
\bottomrule
\end{tabularx}

\textbf{TI sugerido:} $\sim$1\,h práctica plantillas (B) + $\sim$0{,}5\,h lecturas (A)

\subsection{Sesión 3 --- Libros/ebooks I --- Arquitectura y outline}
\textbf{Producto (E3):} Arquitectura de obra + outline con RA

\begin{tabularx}{\textwidth}{@{}lllX@{}}
\toprule
Min & Bloque & Tipo & Actividad \\
\midrule
0:00--0:10 & Recap brief/orquestación & E & De contrato a estructura \\
0:10--0:40 & Arquitectura de obra STEM & E & Capítulos, RA, notación, glosario \\
0:40--1:00 & Demo: brief $\rightarrow$ outline & E & Nivel y verbos de Bloom \\
1:00--1:10 & \emph{Descanso} & --- & --- \\
1:10--2:20 & Taller: arquitectura + outline & P & 6--10 nodos; mapa de símbolos \\
2:20--2:45 & Revisión rápida de outlines & C & Pares: coherencia pedagógica \\
2:45--3:00 & Cierre & C & Avance outline $\rightarrow$ primera sección (C) \\
\bottomrule
\end{tabularx}

\subsection{Sesión 4 --- Libros/ebooks II --- Secciones, ejercicios, consistencia}
\textbf{Producto (E4):} Secciones + artefacto didáctico + consistencia

\begin{tabularx}{\textwidth}{@{}lllX@{}}
\toprule
Min & Bloque & Tipo & Actividad \\
\midrule
0:00--0:10 & Recap outlines & E & Priorizar 1--2 secciones \\
0:10--0:35 & Demo: sección con restricciones & E & Regenerar párrafo débil \\
0:35--0:55 & Figuras, tablas y ejercicios & E & Alineación a objetivos STEM \\
0:55--1:05 & \emph{Descanso} & --- & --- \\
1:05--2:00 & Taller: secciones + ejercicio/figura & P & Generación supervisada \\
2:00--2:35 & Agente de consistencia & P & Notación, unidades, glosario \\
2:35--2:50 & Puesta en común & C & Diff narrado v0$\rightarrow$v1 \\
2:50--3:00 & Cierre & C & Pulir sección (C) + plantilla (B) \\
\bottomrule
\end{tabularx}

\subsection{Sesión 5 --- Presentaciones I --- Objetivos, narrativa, storyboard}
\textbf{Producto (E5):} Storyboard narrativo

\begin{tabularx}{\textwidth}{@{}lllX@{}}
\toprule
Min & Bloque & Tipo & Actividad \\
\midrule
0:00--0:10 & Recap: de texto largo a clase & E & Transferencia \\
0:10--0:40 & Objetivos $\rightarrow$ mensajes $\rightarrow$ arco & E & Carga cognitiva \\
0:40--1:00 & Demo: storyboard & E & 8--12 nodos \\
1:00--1:10 & \emph{Descanso} & --- & --- \\
1:10--2:25 & Taller: storyboard de unidad & P & Mensajes clave + malentendidos \\
2:25--2:50 & Revisión por pares & C & Checklist narrativa \\
2:50--3:00 & Cierre & C & Refinar storyboard (C) \\
\bottomrule
\end{tabularx}

\subsection{Sesión 6 --- Presentaciones II --- Deck, notas, accesibilidad}
\textbf{Producto (E6):} Deck MV + notas + accesibilidad

\begin{tabularx}{\textwidth}{@{}lllX@{}}
\toprule
Min & Bloque & Tipo & Actividad \\
\midrule
0:00--0:10 & Recap storyboards & E & Priorizar 8--12 slides \\
0:10--0:35 & Demo: deck + notas de orador & E & Voz docente; variantes \\
0:35--0:55 & Accesibilidad y densidad visual & E & Contraste, ritmo, notación \\
0:55--1:05 & \emph{Descanso} & --- & --- \\
1:05--2:15 & Taller: deck MV + notas & P & Clase y/o conferencia \\
2:15--2:40 & Mini-revisión accesibilidad & P/C & Checklist rápida \\
2:40--3:00 & Cierre & C & Definir microartefacto final EIC \\
\bottomrule
\end{tabularx}

\subsection{Sesión 7 --- Artículos I --- Estructura, borrador, citas}
\textbf{Producto (E7):} Estructura + borrador con citas a verificar

\begin{tabularx}{\textwidth}{@{}lllX@{}}
\toprule
Min & Bloque & Tipo & Actividad \\
\midrule
0:00--0:15 & Marco de artículo académico & E & IMRyD; audiencia; contribución \\
0:15--0:40 & Demo: borrador asistido + citas & E & Marcas \texttt{[VERIFICAR]} \\
0:40--0:50 & \emph{Descanso corto} & --- & --- \\
0:50--2:15 & Taller: estructura + borrador & P & Micro-sección o guía de lab \\
2:15--2:40 & Lista de verificación & P & Citas/datos \\
2:40--3:00 & Cierre & C & Completar borrador (C) + agent log (D) \\
\bottomrule
\end{tabularx}

\subsection{Sesión 8 --- Artículos II --- Revisión, cierre, showcase}
\textbf{Producto (E8 / acreditación):} Microartefacto + agent log + checklist

\begin{tabularx}{\textwidth}{@{}lllX@{}}
\toprule
Min & Bloque & Tipo & Actividad \\
\midrule
0:00--0:15 & Agentes de revisión & E & Claridad, rigor, sesgos \\
0:15--0:40 & Demo: tres pases + HITL & E & Diff v0$\rightarrow$v1 \\
0:40--0:50 & \emph{Descanso corto} & --- & --- \\
0:50--2:00 & Laboratorio de cierre & P & Integrar pipeline (sesiones 1--7) \\
2:00--2:25 & Agent log + checklist & P & Documentación de acreditación \\
2:25--2:50 & Showcase / galería & C & 3--5 muestras \\
2:50--3:00 & Evaluación del CADi & C & Encuesta / rúbrica (\ceddie) \\
\bottomrule
\end{tabularx}

\section{Parte B --- Trabajo fuera de aula (16\,h)}

\begin{tabularx}{\textwidth}{@{}Xccc@{}}
\toprule
Bloque & Horas & Ventana típica & Apoya \\
\midrule
A. Lecturas y demos guiadas & 4\,h & Entre sesiones 1--2 y refuerzo en 7 & Fundamentos, stack, integridad \\
B. Práctica con stack y plantillas & 4\,h & Entre sesiones 2--6 & Brief, prompt, orquestación \\
C. Avance del microartefacto & 6\,h & Entre sesiones 3--8 & Outline$\rightarrow$secciones / deck / borrador \\
D. Agent log + checklist + reflexión & 2\,h & Sesiones 7--8 y cierre & Acreditación \\
\textbf{Total} & \textbf{16\,h} & A lo largo del CADi & --- \\
\bottomrule
\end{tabularx}

\subsection{Distribución sugerida por intervalo (ejemplo)}

\begin{tabularx}{\textwidth}{@{}Xccccc@{}}
\toprule
Intervalo & A & B & C & D & Suma \\
\midrule
Pre / post sesión 1 & 1{,}0 & --- & --- & --- & 1{,}0 \\
Entre 1 y 2 & 1{,}0 & 0{,}5 & --- & --- & 1{,}5 \\
Entre 2 y 3 & 0{,}5 & 1{,}0 & 0{,}5 & --- & 2{,}0 \\
Entre 3 y 4 & --- & 0{,}5 & 1{,}0 & --- & 1{,}5 \\
Entre 4 y 5 & 0{,}5 & 0{,}5 & 1{,}0 & --- & 2{,}0 \\
Entre 5 y 6 & --- & 0{,}5 & 1{,}0 & --- & 1{,}5 \\
Entre 6 y 7 & 0{,}5 & 0{,}5 & 1{,}0 & --- & 2{,}0 \\
Entre 7 y 8 & 0{,}5 & --- & 1{,}0 & 0{,}5 & 2{,}0 \\
Post sesión 8 & --- & --- & 0{,}5 & 1{,}5 & 2{,}0 \\
\textbf{Total} & \textbf{4{,}0} & \textbf{4{,}0} & \textbf{6{,}0} & \textbf{2{,}0} & \textbf{16{,}0} \\
\bottomrule
\end{tabularx}

\section{Vista rápida}

\begin{tabularx}{\textwidth}{@{}cXl@{}}
\toprule
Sesión & Enfoque & Evidencia \\
\midrule
1 & Fundamentos + integridad EIC & E1 \\
2 & Briefs, prompts, orquestación & E2 \\
3 & Libros I --- arquitectura / outline & E3 \\
4 & Libros II --- secciones / consistencia & E4 \\
5 & Presentaciones I --- narrativa / storyboard & E5 \\
6 & Presentaciones II --- deck / notas / accesibilidad & E6 \\
7 & Artículos I --- estructura / borrador / citas & E7 \\
8 & Artículos II --- revisión / cierre / showcase & E8 \\
\bottomrule
\end{tabularx}

\textbf{Microartefacto (opciones EIC):} capítulo de libro de curso STEM · deck de unidad con notación técnica · sección tipo IMRyD · guía de laboratorio breve.

\textbf{Total:} 40 horas (24 en aula + 16 independientes).
